% Generated by data/prompts/make_app_prompts.py from code/countdown/prompts.py (do not edit by hand: rerun it).
\subsection*{Policy prompt}
\noindent\textbf{User}
\begin{Verbatim}[frame=single, fontsize=\small, breaklines, breakanywhere=false]
Using the numbers 25, 4, 7, 3, create an expression that equals 96. You can use +, -, * and /, and each number must be used exactly once.

Respond in the following format:
<reasoning>
{your search for a way to reach the target}
</reasoning>
<solution>
{the steps, one per line, e.g. 25 * 4 = 100}
</solution>
<answer>
{one expression that uses each number exactly once and equals the target}
</answer>
\end{Verbatim}

\subsection*{Judge prompt}
\noindent\textbf{System}
\begin{Verbatim}[frame=single, fontsize=\small, breaklines, breakanywhere=false]
Please solve the following problem step-by-step, making sure that each step is well-justified.
\end{Verbatim}
\noindent\textbf{User}
\begin{Verbatim}[frame=single, fontsize=\small, breaklines, breakanywhere=false]
Using the numbers 25, 4, 7, 3, create an expression that equals 96. You can use +, -, * and /, and each number must be used exactly once.
\end{Verbatim}
\noindent\textbf{Assistant (scored)}
\begin{Verbatim}[frame=single, fontsize=\small, breaklines, breakanywhere=false]
25 * 4 = 100
100 - 7 = 93
93 + 3 = 96
Answer: 25 * 4 - 7 + 3
\end{Verbatim}

\subsection*{Solution from corrupted reasoning}
The policy prompt above, then the assistant turn starts with:
\noindent\textbf{Assistant (forced start)}
\begin{Verbatim}[frame=single, fontsize=\small, breaklines, breakanywhere=false]
<reasoning>
{kept part of the reasoning, empty for no reasoning}
</reasoning>
<solution>
\end{Verbatim}


\noindent Chat templates: Llama-3.2-3B's template starts every prompt with a system turn holding ``Cutting Knowledge Date: December 2023'' and ``Today Date: 26 Jul 2024'' (the date pinned), followed by the system prompt when there is one. Qwen3-0.6B runs with thinking off: with \texttt{enable\_thinking=False}, its template starts the assistant turn with an empty thinking block, \texttt{<think>} and \texttt{</think>} with a blank line between them.
